\chapter{Twierdzenie Poincarégo w wymiarze dwa i w wysokich wymiarach}
\label{ch:poincare-2-high}
\index{twierdzenie Poincarego@Twierdzenie Poincarégo}

Twierdzenie Poincarégo pyta, kiedy rozmaitość
o typie homotopijnym sfery musi być sferą.
W rozdziale o chirurgii widzieliśmy już,
że \emph{homeomorfizm} i \emph{dyfeomorfizm}
to różne wymagania. Sformułujemy tu
wniosek topologiczny dla powierzchni
i dla gładkich sfer homotopijnych
wymiaru co najmniej pięć.

\begin{definicja}[Sfera homotopijna]
\index{sfera homotopijna@Sfera homotopijna}
\label{def:homotopy-sphere-24}
Zamknięta, spójna rozmaitość $n$-wymiarowa
$\Sigma^n$ jest \emph{sferą homotopijną},
jeśli $\Sigma\simeq S^n$.
Na przykład standardowa $S^n$ spełnia
ten warunek. Definicja nie zawiera
stwierdzenia o dyfeomorfizmie.
\end{definicja}

\section{Powierzchnie: klasyfikacja i konsekwencja}
\label{sec:poincare-2}

\begin{twierdzenie}[Klasyfikacja zamkniętych powierzchni]
\label{thm:classification-surfaces-24}
Każda spójna, zwarta powierzchnia
bez brzegu jest homeomorficzna
dokładnie z jednym modelem z listy:
$S^2$, suma spójna $g\geq1$ torusów
albo suma spójna $h\geq1$ płaszczyzn
rzutowych $\mathbb{RP}^2$.
W pierwszej rodzinie charakterystyka
Eulera wynosi $2-2g$, w drugiej $2-h$;
orientowalność odróżnia rodziny.
\end{twierdzenie}

\begin{proof}
Jedynym wynikiem wejściowym o powierzchniach
jest ich \emph{triangulowalność}: zwarta
topologiczna powierzchnia ma skończoną
triangulację. Dla powierzchni gładkich jest
to twierdzenie przytoczone w
Dodatku~\ref{app:wyniki-zewnetrzne}; wersję
topologiczną udowadnia
\href{https://pi.math.cornell.edu/~hatcher/Papers/TorusTrick.pdf}
{A.\ Hatcher, \emph{The Kirby Torus Trick for Surfaces}}.
Pozostałe kroki rozpiszemy. Oznaczmy przez
$v,e,t$ liczby wierzchołków, krawędzi
i trójkątów wybranej triangulacji $F$.

\emph{Krok 1: charakterystyka i sfera.}
W szkielecie jednowymiarowym wybierzmy
drzewo rozpinające $T$. Jego małe regularne
otoczenie jest dyskiem: zaczynamy od dysku
przy jednym wierzchołku, a każda następna
krawędź drzewa dokleja prostokątny pas
do dotychczasowego dysku wzdłuż jednego łuku.
W szczególności $T$ nie rozcina $F$;
ścieżkę przecinającą $T$ można poprowadzić
wzdłuż brzegu tego otoczenia.

Niech $D$ będzie grafem dualnym z wierzchołkiem
w każdym trójkącie i krawędzią przecinającą
każdą krawędź triangulacji \emph{spoza} $T$.
Jest spójny: przechodząc między trójkątami
w $F\setminus T$, możemy omijać krawędzie
$T$, więc kolejne przejścia dają drogę w $D$.
Graf $T$ ma $v$ wierzchołków i $v-1$ krawędzi,
a $D$ ma $t$ wierzchołków i $e-(v-1)$ krawędzi.
Stąd, pisząc $b_1(D)=|E(D)|-|V(D)|+1$,
otrzymujemy
\[
 \chi(F)=v-e+t=1+\chi(D)=2-b_1(D)\leq 2.
\]
Jeżeli $\chi(F)=2$, to $D$ jest drzewem.
Pogrubienia obu drzew są dyskami;
prostokąty wzdłuż krawędzi triangulacji
spoza $T$ wypełniają obszar między ich
brzegami. Dają zatem rozkład $F$ na dwa
dyski sklejone całymi brzegami, czyli
$F\cong S^2$. Odwrotna implikacja wynika
z $\chi(S^2)=2$.

Jeżeli $\chi(F)<2$, graf $D$ zawiera prosty
cykl $\gamma$. Jest on \emph{nierozdzielający}:
$F\setminus D$ jest spójne, bo każdy punkt
można połączyć w swoim trójkącie z wierzchołkiem
triangulacji, a wszystkie te wierzchołki łączy
$T\subset F\setminus D$. Zatem również
$F\setminus\gamma$ jest spójne.

\emph{Krok 2: cięcie i zmiana $\chi$.}
Prosta zamknięta krzywa ma regularne otoczenie
będące pierścieniem albo wstęgą Möbiusa.
Rozcięcie wzdłuż krzywej nie zmienia
charakterystyki Eulera: zarówno jej
otoczenie, jak i okręgi wspólnego brzegu
mają $\chi=0$, więc wzór na $\chi$ sumy
powierzchni daje $\chi(F)=\chi(F\setminus
\operatorname{int}N(\gamma))$.
Zamknięcie każdej nowej składowej brzegu
dyskiem zwiększa $\chi$ o jeden.
W przypadku krzywej dwustronnej są dwa
okręgi brzegowe; w przypadku jednostronnej
jest jeden. Krzywa jednostronna nie rozdziela
$F$, gdyż dopełnienie jej wstęgi Möbiusa
ma tylko jeden okrąg brzegowy.

\emph{Krok 3: powierzchnie orientowalne.}
Indukujemy po $r=2-\chi(F)$. Dla $r=0$
mamy sferę z Kroku 1. Jeśli $F$ jest
orientowalna i $r>0$, wybieramy powyższy
nierozdzielający cykl $\gamma$. Jest dwustronny.
Po rozcięciu i zamknięciu obu brzegów
dostajemy spójną orientowalną powierzchnię
$F'$ o $\chi(F')=\chi(F)+2$.
Z założenia indukcyjnego $F'\cong\#^gT^2$
dla pewnego $g\geq0$. Aby odzyskać $F$,
usuwamy z $F'$ wnętrza dwóch dysków i
łączymy ich brzegi pierścieniem.
Zgodność orientacji wymusza zwykły
uchwyt; jest to dokładnie operacja
$F\cong F'\#T^2$ opisana w
Przykładzie~\ref{prz:sfera-torus-chirurgia-18}.
Zatem $F\cong\#^{g+1}T^2$.

\emph{Krok 4: powierzchnie nieorientowalne.}
Istnieje na $F$ zamknięta droga odwracająca
orientację. Po ustawieniu jej w położeniu
ogólnym wybieramy taką drogę z najmniejszą
liczbą samoprzecięć. W każdym podwójnym
punkcie wygładzenie rozcina drogę na dwie;
iloczyn ich znaków zmiany orientacji jest
równy $-1$, więc jedna z nich także odwraca
orientację. Wybierając ją, zmniejszamy
liczbę samoprzecięć. W końcu otrzymujemy
prostą krzywą jednostronną $\delta$.
Po rozcięciu wzdłuż $\delta$ i zamknięciu
jedynego brzegu dyskiem powstaje spójna
powierzchnia $F'$ z $\chi(F')=\chi(F)+1$.
Odwrotnie, $F\cong F'\#\mathbb{RP}^2$,
bo zastępujemy ten dysk wstęgą Möbiusa.

Jeżeli $F'$ jest nieorientowalna, indukcja
daje $F'\cong\#^h\mathbb{RP}^2$, a więc
$F\cong\#^{h+1}\mathbb{RP}^2$.
Jeżeli $F'$ jest orientowalna, to
$F\cong(\#^gT^2)\#\mathbb{RP}^2$.
Przy $g=0$ dostajemy $\mathbb{RP}^2$.
Przy $g>0$ stosujemy lokalne przesunięcie
wstęg
\begin{equation}\label{eq:crosscap-handle-24}
 T^2\#\mathbb{RP}^2
 \cong\mathbb{RP}^2\#\mathbb{RP}^2
       \#\mathbb{RP}^2.
\end{equation}
Aby sprawdzić tę relację, przedstawmy
lewą stronę jako dysk z trzema pasami:
dwa nieskręcone pasy torusa mają końce
ułożone naprzemiennie, trzeci pas ma
półobrót i daje płaszczyznę rzutową.
Brzeg wielokąta opisuje słowo
$aba^{-1}b^{-1}cc$. Użyjemy elementarnego
przekształcenia
\[
 xUxV\ \sim\ U^{-1}xxV.
\]
Obie litery $x$ oznaczają boki sklejane
w tym samym kierunku. Rozcinamy wielokąt
wzdłuż łuku łączącego te boki i sklejamy
go po drugiej stronie pary: boki $x$
stają obok siebie, a odcinek brzegu $U$
jest teraz obchodzony w odwrotnym kierunku.
Iloraz powierzchni się nie zmienia.
Wolno też cyklicznie przesuwać słowo
brzegowe i zastępować nazwę boku
jej odwrotnością. Stosując tę regułę
cztery razy, dostajemy jawnie
\[
\begin{aligned}
 aba^{-1}b^{-1}cc
 &\sim aba^{-1}cbc
  \sim abbc^{-1}ac\\
 &\sim aacb^{-1}b^{-1}c
  \sim aabbcc.
\end{aligned}
\]
Końcowe trzy pary boków są skręcone,
więc wyznaczają sumę spójną trzech
$\mathbb{RP}^2$, co dowodzi~\eqref{eq:crosscap-handle-24}.
Ten sam rachunek słów zapisuje
\href{https://pnp.mathematik.uni-stuttgart.de/igt/eiserm/lehre/Topologie/Topologie-K-2x2.pdf}
{M.\ Eisermann, \emph{Topologie}, twierdzenie K3f (8)};
diagram geometryczny podaje
\href{https://www.math.uchicago.edu/~may/VIGRE/VIGRE2008/REUPapers/Huang.pdf}
{J.\ Huang, \emph{Classification of Surfaces}, rys.~13}.
Stosując~\eqref{eq:crosscap-handle-24}
$g$ razy, otrzymujemy sumę spójną
$(2g+1)$ płaszczyzn rzutowych.

Pozostaje jednoznaczność. Dodawanie torusa
zmniejsza $\chi$ o dwa, a dodawanie
$\mathbb{RP}^2$ o jeden, gdyż
$\chi(A\#B)=\chi(A)+\chi(B)-2$.
Tak dostajemy podane w tezie wzory.
Obie rodziny rozróżnia orientowalność,
a wewnątrz każdej rodziny wartość $\chi$
wyznacza liczbę składników. Charakterystyka
jest niezmiennikiem topologicznym, ponieważ
równa się naprzemiennej sumie rang homologii
singularnych z Dodatku~\ref{app:homologia-singularna-euler}.
\end{proof}

W powyższym dowodzie zewnętrzne jest tylko
twierdzenie o triangulacji powierzchni.
Krótki wariant argumentu dla powierzchni
orientowalnych podaje
\href{https://academicweb.nd.edu/~andyp/notes/ClassificationSurfaces.pdf}
{A.\ Putman, \emph{A Quick Proof of the Classification of Surfaces}};
pełniejszą wersję obu przypadków zawierają
\href{https://academicweb.nd.edu/~andyp/notes/ClassificationSurfacesLong.pdf}
{jego notatki, §5--6}.

\begin{lemat}[Pierwsza homologia modeli]
\label{lem:H1-surfaces-24}
Dla $g,h\geq1$ mamy
\[
 H_1(\#^g T^2;\Z)\cong\Z^{2g},
 \qquad
 H_1(\#^h\mathbb{RP}^2;\Z)
     \cong\Z^{h-1}\oplus\Z/2.
\]
\end{lemat}
\begin{proof}
Standardowy wielokąt orientowalnej
powierzchni ma kolejno oznaczone boki
\[
a_1,b_1,a_1^{-1},b_1^{-1},\ldots,
a_g,b_g,a_g^{-1},b_g^{-1}.
\]
Po sklejeniu dostajemy jedną komórkę
zerową, $2g$ komórek jednowymiarowych
i jedną dwuwymiarową. Operator brzegu
$\partial_1$ jest zerowy, bo oba końce
każdej krawędzi to ten sam wierzchołek.
W brzegowym słowie komórki dwuwymiarowej
każda krawędź pojawia się raz z plusem
i raz z minusem, więc $\partial_2=0$.
Zatem $H_1=\ker\partial_1/
\operatorname{im}\partial_2=\Z^{2g}$.

Nieorientowalny wielokąt ma słowo
$a_1a_1\cdots a_ha_h$. Teraz
$C_1\cong\Z^h$, $\partial_1=0$,
a $\partial_2(1)=2(a_1+\cdots+a_h)$.
Wektor $(1,\ldots,1)$ jest prymitywny
w $\Z^h$, więc można go uzupełnić do
bazy całkowitej. W tej bazie iloraz
przez podgrupę generowaną przez jego
podwojenie jest
$\Z^{h-1}\oplus\Z/2$.
\end{proof}

\begin{twierdzenie}[Twierdzenie Poincarégo dla powierzchni]
\label{thm:poincare-2-24}
Każda spójna, zamknięta powierzchnia
homotopijnie równoważna $S^2$
jest homeomorficzna z $S^2$.
Równoważnie: jedyną spójną, zamkniętą
i po prostu spójną powierzchnią
jest $S^2$.
\end{twierdzenie}
\begin{proof}
Równoważność homotopijna zachowuje
homologię, więc $H_1(\Sigma;\Z)=
H_1(S^2;\Z)=0$.
Twierdzenie~\ref{thm:classification-surfaces-24}
ogranicza $\Sigma$ do trzech rodzin.
Lemat~\ref{lem:H1-surfaces-24}
wyklucza wszystkie modele o $g\geq1$
i $h\geq1$, ponieważ każda z ich
grup $H_1$ jest niezerowa.
Pozostaje tylko $S^2$.
Jeśli założymy zamiast równoważności
homotopijnej prostą spójność, to
$H_1(\Sigma;\Z)$ jest abelianizacją
$\pi_1(\Sigma)$, więc ten sam argument
daje identyczny wniosek.
\end{proof}

\section{Usunięcie dwóch dysków ze sfery homotopijnej}
\label{sec:poincare-high}

Następny argument dotyczy \emph{gładkiej}
sfery homotopijnej $\Sigma^n$.
Wybierzmy dwa rozłączne gładko
osadzone zamknięte dyski $D_-^n,D_+^n$
i połóżmy
\[
 W=\Sigma\setminus
   \bigl(\operatorname{int}D_-^n
          \sqcup\operatorname{int}D_+^n\bigr),
 \qquad
 \partial W=S_-^{n-1}\sqcup S_+^{n-1}.
\]
Przestrzeń $W$ jest kobordyzmem,
którego dwiema składowymi brzegu
są sfery standardowe.

\begin{lemat}[Dlaczego to $h$-kobordyzm?]
\label{lem:punctured-sphere-hc-24}
Jeżeli $n\geq3$ i $\Sigma^n\simeq S^n$,
to obie inkluzje
$S_\pm^{n-1}\hookrightarrow W$
są równoważnościami homotopijnymi.
\end{lemat}
\begin{proof}
Przez wycinanie para $(\Sigma,W)$
ma tę samą homologię względną
co suma dwóch par $(D^n,S^{n-1})$.
Jej jedyną niezerową grupą względną
w dodatnim stopniu jest zatem
$H_n(\Sigma,W)\cong\Z\oplus\Z$.
Fragment długiego ciągu ma postać
\[
 0\longrightarrow H_n(\Sigma)\cong\Z
 \xrightarrow{\ (1,1)\ }
 \Z\oplus\Z
 \xrightarrow{\ \partial\ }
 H_{n-1}(W)\longrightarrow0.
\]
Znaki obu współrzędnych zależą od
orientacji dysków; wybieramy je tak,
aby pierwsza strzałka była diagonalna.
Stąd $H_{n-1}(W)\cong\Z$, a obrazy
obu wektorów bazowych przez $\partial$
są jej generatorami z przeciwnymi znakami.
Są to klasy fundamentalne obu składowych
brzegu. Pozostałe grupy zredukowane
$H_i(W)$ znikają, jak dla $S^{n-1}$.

Przez twierdzenie van Kampena,
usunięcie wnętrza każdego $n$-dysku
nie zmienia $\pi_1$ dla $n\geq3$:
część wspólna w odpowiednim otwartym
rozkładzie deformuje się na
$S^{n-1}$ i jest po prostu spójna.
Zatem $\pi_1(W)\cong\pi_1(\Sigma)=1$.
Obie sfery brzegowe też są po prostu
spójne. Ich inkluzje są izomorfizmami
homologii, więc homologiczną wersją
twierdzenia Whiteheada
(Wniosek~\ref{cor:whitehead-homologiczny})
są równoważnościami homotopijnymi.
\end{proof}

\begin{figure}[htbp]
\centering
\begin{tikzpicture}[>=Latex,font=\small,line cap=round]
 \shade[inner color=white,outer color=manifoldblue!43]
      (0,0) circle (1.55);
 \draw[manifoldblue!75!black,very thick]
      (0,0) circle (1.55);
 \fill[white] (-1.02,.70) circle (.27);
 \fill[white] (1.02,-.70) circle (.27);
 \draw[manifoldgold!80!black,thick]
      (-1.02,.70) circle (.27)
      (1.02,-.70) circle (.27);
 \node[anchor=east,align=right] at (-1.42,1.15)
      {$D_-^n$\\usunięty};
 \node[anchor=west,align=left] at (1.42,-1.13)
      {$D_+^n$\\usunięty};
 \draw[->,very thick,manifoldblue!70!black]
      (2.05,0)--(3.15,0);
 \draw[manifoldblue!75!black,very thick,
       fill=manifoldblue!12]
      (3.52,-1.07) rectangle (6.82,1.07);
 \draw[manifoldgold!85!black,very thick]
      (3.52,-1.07)--(3.52,1.07)
      (6.82,-1.07)--(6.82,1.07);
 \node at (5.17,.34) {$W$};
 \node at (3.52,-1.40) {$S_-^{n-1}$};
 \node at (6.82,-1.40) {$S_+^{n-1}$};
\end{tikzpicture}
\caption{Schemat wycięcia dwóch dysków z $\Sigma^n$.
Prostokąt przedstawia kobordyzm po zastosowaniu
twierdzenia o $h$-kobordyzmie; dla $n\geq6$
jest on gładkim iloczynem. Rysunek jest
schematyczny, nie jest osadzeniem $\Sigma^n$
w płaszczyźnie.}
\label{fig:poincare-two-disks-24}
\end{figure}

\begin{twierdzenie}[Poincaré dla gładkich
sfer homotopijnych, $n\geq6$]
\label{thm:poincare-high-24}
Jeżeli $\Sigma^n$ jest gładką sferą
homotopijną i $n\geq6$, to
$\Sigma$ jest homeomorficzna z $S^n$.
\end{twierdzenie}
\begin{proof}
Lemat~\ref{lem:punctured-sphere-hc-24}
mówi, że $(W;S_-^{n-1},S_+^{n-1})$
jest po prostu spójnym $h$-kobordyzmem.
Jego wymiar wynosi $n$, zatem $n\geq6$
pozwala użyć gładkiego twierdzenia
Smale'a o $h$-kobordyzmie
(Twierdzenie~\ref{thm:smale-hc-19}).
Otrzymujemy dyfeomorfizm
$W\cong S^{n-1}\times[0,1]$
ustalony na dolnym brzegu.
Po ponownym doklejeniu dwóch dysków
powstaje przestrzeń
$D^n\cup_{\varphi}D^n$
dla pewnego dyfeomorfizmu
$\varphi:S^{n-1}\to S^{n-1}$.
Każdy taki $\varphi$ ma \emph{ciągłe}
przedłużenie na dysk dane wzorem
$r u\mapsto r\varphi(u)$ dla
$0\leq r\leq1$, $u\in S^{n-1}$;
w środku kładziemy $0\mapsto0$.
Wzór odwrotny pochodzi od
$\varphi^{-1}$, więc jest to
homeomorfizm dysków. Zmieniając
nim parametryzację jednego dysku,
uzyskujemy homeomorfizm
$D^n\cup_\varphi D^n\cong
D^n\cup_{\id}D^n\cong S^n$.
Stożkowe przedłużenie nie musi
być gładkie w środku dysku;
dlatego dowód daje homeomorfizm,
a nie dyfeomorfizm.
\end{proof}

\section{Osobny przypadek wymiaru pięć}
\label{sec:poincare-five}

Dla $n=5$ kobordyzm $W$ z poprzedniej
sekcji ma wymiar $5$, podczas gdy
Twierdzenie~\ref{thm:smale-hc-19}
wymaga wymiaru kobordyzmu co najmniej $6$.
Nie wolno zastosować go tutaj przez
samą zmianę numeru wymiaru.
Można jednak bez klasyfikacji rozpoznać
topologię sfery po wyjęciu \emph{jednego}
dysku; ujawnia to dokładnie brakujący krok.

\begin{lemat}[Co daje wyjęcie jednego dysku]
\label{lem:punctured-five-24}
Jeżeli $\Sigma^5\simeq S^5$ jest gładka
i $D^5\subset\Sigma$ jest gładko osadzonym
dyskiem, to
$C=\Sigma\setminus\operatorname{int}D^5$
jest zwartą gładką rozmaitością
\emph{ściągalną} o brzegu dyfeomorficznym
z $S^4$.
\end{lemat}
\begin{proof}
Gładka tuba brzegu dysku daje rozkład
$\Sigma=C\cup D^5$ z częścią wspólną
o typie homotopijnym $S^4$.
W ciągu Mayera--Vietorisa homomorfizm
$H_5(\Sigma)\cong\Z\to H_4(S^4)\cong\Z$
jest izomorfizmem: wysyła klasę
fundamentalną na klasę brzegu dysku.
Z dokładności dostajemy $H_4(C)=0$,
a wszystkie pozostałe dodatnie grupy
homologii $C$ też znikają, ponieważ
$\Sigma$ ma homologię $S^5$.
Twierdzenie van Kampena daje
$\pi_1(C)=\pi_1(\Sigma)=1$,
gdyż $D^5$ i $S^4$ są po prostu spójne.
Rozmaitość ma typ kompleksu CW;
homologiczna wersja twierdzenia
Whiteheada zastosowana do
punktu $*\to C$ daje ściągalność.
Brzeg $C$ jest dokładnie
$\partial D^5\cong S^4$.
\end{proof}

Ściągalność $C$ nie daje jeszcze
dyfeomorfizmu $C\cong D^5$.
Takie wnioskowanie wymagałoby
oddzielnego twierdzenia o gładkich
wypełnieniach $S^4$. W następnym
argumencie korzystamy z klasyfikacji
zamkniętych $5$-rozmaitości Barden.
Posłużymy się następującym dokładnym
wynikiem zewnętrznym.

\begin{twierdzenie}[Barden, twierdzenie 2.2;
 wynik zewnętrzny]
\label{thm:Barden-5-24}
Niech $P,Q$ będą zamkniętymi, zorientowanymi,
gładkimi, po prostu spójnymi
$5$-rozmaitościami. Jeśli izomorfizm
$H_2(P;\Z)\to H_2(Q;\Z)$ zachowuje
formę wiązania na części torsyjnej
i klasy $w_2$, to pochodzi od
dyfeomorfizmu zachowującego orientację.
\end{twierdzenie}

Jest to
\href{https://www.sas.rochester.edu/mth/sites/doug-ravenel/otherpapers/barden.pdf}
{D.\ Barden, \emph{Simply Connected Five-Manifolds},
Ann. of Math. 82 (1965), twierdzenie 2.2}.
Jego dowód konstruuje odpowiedni
$6$-wymiarowy kobordyzm i jest
oddzielnym głębokim argumentem.
Warto zobaczyć, gdzie pojawia się
poprzednio brakujący wymiar.
Barden buduje kobordyzm $V^6$ między
porównywanymi $5$-rozmaitościami,
kontrolując kolejno $\pi_1(V)$,
$H_2(V)$ i klasy obramowań uchwytów.
Po dodatkowych modyfikacjach obie
inkluzje brzegów stają się
równoważnościami homotopijnymi;
otrzymuje się więc $h$-kobordyzm
\emph{wymiaru sześć}.
Teraz Twierdzenie~\ref{thm:smale-hc-19}
można zastosować zgodnie z jego
założeniami. Konstrukcja i kontrola
uchwytów są treścią §3--5 pracy
Bardena, a nie konsekwencją samego
Lematu~\ref{lem:punctured-five-24}.

\begin{wniosek}[Poincaré dla gładkich
sfer homotopijnych wymiaru pięć]
\label{cor:poincare-five-24}
Każda gładka $\Sigma^5\simeq S^5$
jest dyfeomorficzna z $S^5$.
\end{wniosek}
\begin{proof}
Równoważność homotopijna daje
$\pi_1(\Sigma)=1$ oraz
$H_2(\Sigma;\Z)=H_2(S^5;\Z)=0$.
Z twierdzenia o współczynnikach
uniwersalnych wynika również
$H^2(\Sigma;\Z/2)=0$, bo
$H_1(\Sigma)=H_2(\Sigma)=0$.
Zatem $w_2(\Sigma)=0$, jak dla $S^5$.
Jedyny izomorfizm grup zerowych
$H_2(\Sigma)\to H_2(S^5)$
zachowuje trywialną formę na grupie zerowej
i obie zerowe klasy $w_2$.
Twierdzenie~\ref{thm:Barden-5-24}
daje dyfeomorfizm.
\end{proof}

Wyniki tej sekcji i
Twierdzenie~\ref{thm:poincare-high-24}
razem dowodzą, że \emph{każda
gładka} sfera homotopijna wymiaru
co najmniej pięć jest
\emph{homeomorficzna} ze sferą
standardową. W wymiarze pięć
otrzymaliśmy nawet dyfeomorfizm;
w wymiarze siedem ten mocniejszy
wniosek jest fałszywy, co pokażemy
w następnym rozdziale.

\clearpage
